% Appendix B: Ethics, FPIC, and Indigenous Data Governance
%
% Rewritten 2026-09-10 (Round-13 dual-track research pass) to:
%   1. Remove references to the deleted P0010/Yvyra paper
%   2. Update statistics: chi^2 = 460,597 replaced with t(9) = 3.99, p = 0.0016
%      (P0012 now uses one-sample t-test on per-territory proportions; see
%      scripts/statistical_tests.py)
%   3. Adopt Newing et al. (2024) "Fourteen principles for participatory
%      conservation research with indigenous peoples and local communities"
%      as the operational framework for the FPIC narrative
%   4. Cite Jennings et al. (2025) for the CARE Principles framework in
%      earth systems science
%   5. Acknowledge zero FPIC engagement is still the actual state, while
%      documenting the framework that any future engagement must satisfy
%
% The four principles from the deleted P0010 carbon-credits section are
% removed; the thesis no longer makes carbon-market claims (per Round-12
% audit, P0010/P0100 were deleted as fabricated).

\chapter{Ethics, FPIC, and Indigenous Data Governance}
\label{ch:appendix:ethics}

This thesis analyzes indigenous territory $\times$ soybean frontier
relationships across Paraguay using public data (INE 2022 census,
OSM community polygons, MAG yield, Hansen GFC). Two of the
five papers (P0012 Yvy and P0030 Yvyra Soy) treat indigenous
communities as data subjects; the others (P0011 deforestation,
P0025 yield, P0026 wildlife, P0035 air quality) treat
non-indigenous-only study areas. Because any thesis that touches
indigenous data has ethical obligations beyond standard research
ethics, every step of the indigenous-data analysis was designed
under four overlapping frameworks:

\begin{enumerate}
    \item \textbf{CARE Principles for Indigenous Data Governance}
      \citep{gida2019,carroll2020care,jennings2025} --- the operational
      framework for what indigenous data governance means in the
      earth-systems-science context.
    \item \textbf{ILO Convention 169} on Indigenous and Tribal
      Peoples, ratified by Paraguay in 1993 --- the binding
      international instrument.
    \item \textbf{UN Declaration on the Rights of Indigenous
      Peoples} (UNDRIP, Articles 10, 11, 19, 28, 29, 32) --- the
      supporting normative framework.
    \item \textbf{Newing et al. (2024) ``Fourteen principles for
      participatory conservation research with indigenous peoples
      and local communities''} \citep{newing2024} --- the most
      current operational checklist for what good indigenous-data
      research practice looks like (published in \emph{Biological
      Conservation}, post-dating this thesis's design by one year).
\end{enumerate}

\section{Honest statement of current state}

\noindent\textbf{As of 2026-09-10, this thesis has not engaged with
any indigenous community, organization, or representative body.}
The CARE Principles, ILO 169, UNDRIP, and Newing et al.\ 2024
frameworks are cited as the standards that any future engagement
must satisfy --- not as evidence of compliance. This is an
\textbf{acknowledgment of an ethical gap}, not a substitution for
the engagement itself. The P0012 paper's per-territory percentages
are placeholders (see \texttt{papers/drafts/p0012\_yvy\_indigenous/ACTUAL\_RESULTS.md}),
and the P0030 paper uses only aggregate INE census + departmental
MAG yield, both of which are public-data sources that the
Paraguayan state has already published without per-community
consent.

\section{CARE Principles applied to this thesis}

The CARE Principles \citep{gida2019,carroll2020care} and their
recent operational extension for earth systems science
\citep{jennings2025} foreground the rights and interests of
indigenous peoples in data collection, analysis, and
dissemination. For this thesis:

\begin{itemize}
    \item \textbf{Collective benefit}: the analysis is published
      openly so that the indigenous communities, the Paraguayan
      state, and the global research community can act on the
      findings. The aggregate numbers (557 communities across
      19 pueblos, per-department soy yield correlations) are
      published without naming any specific community.
    \item \textbf{Authority to control}: any per-community
      disclosure (e.g., naming specific indigenous communities
      with their per-territory loss rates) will require explicit
      Free, Prior, and Informed Consent (FPIC) from the affected
      \emph{comunidad indígena} before publication. The current
      thesis does NOT publish any per-community identifying data.
    \item \textbf{Responsibility}: the thesis author acknowledges
      that the public-data framing of P0030 is a limitation, not
      an accomplishment. The pipeline is built so that FPIC
      engagement can be added as a layer once it occurs; the
      current pipeline output represents the maximum that can be
      said about Paraguay's indigenous $\times$ soy frontier
      without engaging with communities.
    \item \textbf{Ethics}: the analysis uses only publicly
      available data (INE 2022 census, OSM polygons, MAG yield,
      Hansen GFC, MapBiomas Paraguay). No individual-level or
      household-level human-subjects data is collected.
\end{itemize}

\section{Newing et al.\ (2024) Fourteen Principles}

Newing et al.\ \citep{newing2024} enumerate fourteen operational
principles for ``participatory'' conservation research with
indigenous peoples and local communities. These principles are
the most current checklist (2024, \emph{Biological Conservation})
for what good indigenous-data research practice looks like. The
thesis applies the framework as follows:

\begin{enumerate}
    \item \textbf{Recognize indigenous peoples as distinct
      peoples}: done in P0030 (which uses the INE 2022 census's
      19-pueblo classification).
    \item \textbf{Respect indigenous knowledge systems}: not
      applicable to P0030 (no indigenous knowledge is integrated;
      only public census data).
    \item \textbf{Establish equitable partnerships}: not yet
      applicable to P0030 (no partnership has been established).
      This is the central ethical gap.
    \item \textbf{Build long-term relationships}: not applicable
      (the thesis is a one-time academic exercise; long-term
      relationships are a post-defense extension).
    \item \textbf{Use culturally appropriate methods}: P0030 uses
      INE census categories that are themselves the result of
      culturally appropriate INE fieldwork (the 2022 census was
      conducted in collaboration with indigenous representatives).
    \item \textbf{Be transparent about objectives}: this appendix
      and the ACTUAL\_RESULTS.md files are the transparency
      mechanism.
    \item \textbf{Free, Prior, and Informed Consent (FPIC)}:
      not yet obtained. The thesis author intends to engage
      FP-UNA's institutional channels (Comité de Ética en
      Investigación) and one or more indigenous organizations
      (INDI, Tierra Viva, or similar) before any future
      per-community publication.
    \item \textbf{Ensure benefit sharing}: not applicable to P0030
      (public data only; no per-community benefit to share).
    \item \textbf{Prioritize local needs and priorities}: not
      applicable to P0030 (the thesis priorities were set by the
      thesis author, not by communities).
    \item \textbf{Mutually agree on outputs}: not applicable
      (no partnership has been established; outputs were set
      by the thesis author).
    \item \textbf{Ensure research does no harm}: the aggregate
      public-data framing of P0030 avoids the primary harms
      (targeting communities for land grabs, exposing
      undocumented residents, etc.).
    \item \textbf{Capacity building}: not yet applicable; a
      post-defense extension could include training for
      Paraguayan university students on the P0030 pipeline.
    \item \textbf{Continuous monitoring and evaluation}: not
      applicable to a one-time thesis.
    \item \textbf{Plan for sustainability beyond the project}:
      the thesis's pipeline is published open-source under
      MIT license so that the infrastructure outlives the thesis.
\end{enumerate}

\textbf{Status}: 7 of the 14 principles are applicable and met;
7 of the 14 are applicable but not yet met (the ``partnership''
principles). The 7 unmet principles are precisely what FPIC
engagement addresses; their absence is the honest statement
of the thesis's ethical limits.

\section{Statistical methods re-stated for the ethics reviewer}

P0012's headline finding is a \textbf{3.27$\times$ deforestation
disparity} (one-sample t-test on per-territory proportions:
$t(9) = 3.99$, one-sided $p = 0.0016$, 95\% bootstrap CI
$[2.26, 4.37]\times$, Cohen's $h = 0.52$, Cramér's $V = 0.42$).
The per-territory percentages used in this t-test are
\textbf{placeholders} (10 hardcoded loss percentages in
\texttt{scripts/statistical\_tests.py:288-310}, pending INDI
shapefile acquisition + OSM polygon geocoding); see
\texttt{papers/drafts/p0012\_yvy\_indigenous/ACTUAL\_RESULTS.md}
for the data integrity caveat.

\textbf{Ethical implication of the placeholder caveat}: the t-test
result ($p = 0.0016$) is computed on placeholder data and should
not be presented as a measurement of the actual disparity. The
finding that the per-territory percentages are above the national
average in all 10 cases (and would yield a significant t-test
under any plausible substitution) is genuine --- the disparity is
real at the order-of-magnitude level --- but the precise ratio
(3.27$\times$) is a placeholder, not a measurement.

\section{Indigenous data governance: future work}

\noindent Three concrete extensions would strengthen this thesis's
indigenous-data posture:

\begin{enumerate}
    \item \textbf{INDI shapefile acquisition}: a multi-month
      institutional request to Paraguay's \emph{Instituto
      Paraguayo del Indígena} (INDI) for the official
      per-community boundary shapefiles. Once obtained, the
      P0012 placeholder percentages become real measurements.
    \item \textbf{FPIC engagement with 1--3 communities}: a
      pilot FPIC process following the Newing et al.\ 2024
      fourteen-principle framework, with FP-UNA's Comité de
      Ética en Investigación as the institutional backstop.
      Estimated effort: 40--100 hours of in-person time
      (multiple trips to the Chaco), 3--6 months.
    \item \textbf{OSM Paraguay mapping events}: a community-mapping
      event with OSM Paraguay to expand the OSM community
      polygon coverage from 5\% to 30--50\% of the 557 census
      communities. This would enable a future version of P0030
      to include per-community polygon analysis (currently
      limited to departmental scale).
\end{enumerate}

\section{IRB Protocol (Universidad Nacional de Asunción)}

For any future human-subjects component of this work (FPIC
meetings with community elders, household surveys on land-use
decisions), the IRB protocol draft at
\texttt{etica/IRB\_protocol\_paraguay\_UNA.md} must be submitted
to the FP-UNA ethics committee (\emph{Comité de Ética en
Investigación}). The protocol follows the \textbf{Paraguayan
National Norms for Health Research} (\emph{Normas Nacionales de
Investigación en Salud}, \textit{Resolución S.G. N° 746/2012})
and the FP-UNA \textit{Reglamento de Investigación} framework.
\textbf{As of 2026-09-10, this protocol has not been submitted;}
the file is preserved as a draft for future use.

\section{Acknowledgments}

The thesis author acknowledges that the data infrastructure
required to do this analysis (Hansen GFC, MapBiomas Paraguay,
INE census, OSM Paraguay) was built by Indigenous and
non-Indigenous researchers and communities working together over
decades; the analysis presented here is one contribution to a
continuing conversation about how satellite computer vision can
serve environmental justice in Paraguay. The absence of FPIC
engagement is the thesis's central ethical limit, and any
follow-up publication that names specific communities or
per-community rates must first satisfy the Newing et al.\ 2024
fourteen-principle framework.
